\section{Frequency tools for both branches}\label{sec:prime-estimates}

The shared comparison has left two independent products of large
prime labels. We first record its common Fourier normalization, then
the prime estimates used by the two frequency branches.

\subsection{The common Fourier reduction}

Dyadically localize the products in
\eqref{eq:gr-comparison-term} to \(D_i\asymp U_i\). Their ranges give
\[
 U_i\ge\exp(L^{1/3}),\qquad
 U_i\le\exp(O_m((\log L)L^{2/5})),\qquad Y=xU_1U_2.
\]
After inserting endpoint dyadic cutoffs before Mellin inversion, the
ordinary comparison has the form
\begin{equation}\label{eq:gr-fouriercomparison}
 \frac1Y\int_0^1\widehat F(\theta)\widehat G(-\theta)
                           m(k\theta)\,d\theta,\qquad
 m(\theta)=\E b_1(D_1)b_2(D_2)\ee{\theta D_1D_2}.
\end{equation}
The endpoint reflection for sums changes only signs and fixed character
factors. External Mellin transforms have fixed logarithmic costs.
The centered endpoint has the following common form.
Let \(\xi\) be either centered function in \textup{(C3)}, at a
fixed depth which will be chosen after the required precision.
For a comparison term, define
\begin{equation}
 F(n)=\psi(n/Y)n^{i\sigma}\xi(n)W_{\mathbf t}^b(n),
 \quad Y=x^{1+o(1)},\quad |\sigma|\le L^C,             \label{eq:gr-energy-F}
\end{equation}
where the support of \(\psi\) is a fixed compact positive interval.
All its logarithmic derivatives are bounded by fixed powers of \(L\).
The coefficient \(b\) ignores a selected small-bank label, whose count
is exactly one.  This holds for every comparison above.
Reflection changes only fixed character factors and additive signs,
so we state the estimate at a positive endpoint.


The other endpoint satisfies \(\sum_n|G(n)|^2\ll YL^C\), by
bounded invariant endpoints and the pointwise list bound. No
multiplicativity of the list coefficients is required.

By \eqref{eq:gr-productlaw} and
\eqref{eq:gr-integer-bilinear}, a rational approximation
\(|\theta-a/r|\le r^{-2}\) gives
\begin{equation}\label{eq:gr-multiplier}
 |m(\theta)|^2\ll L^{C_{\rm prod}}
 \left(r^{-1}+U_1^{-1}+\frac{\log(2r)}{U_2}
                    +\frac{r\log(2r)}{U_1U_2}\right).
\end{equation}
The two branches apply this bound at \(k=h\) and \(k=x\),
respectively; their resulting physical major windows differ.

\subsection{Sharp twisted prime sums}


The sharp estimate uses the zero-free input at conductors polynomial
in \(x\). The reflected branch needs that range; the fixed-shift
branch uses the smooth version below at polynomial-logarithmic
conductors and polynomial height.

\begin{lemma}[Prime factor estimate]\label{lem:gr-primes}
Assume the fixed all-conductor zero-free half-plane stated in
Section~\ref{sec:inputs}.
Fix \(\sigma_1\in(11/12,1)\) and \(a,B,B_1>0\).
For \(q\le x^B\), every character \(\chi\bmod q\), and
\(\exp(L^a)\le P\le x^{B_1}\), uniformly on subintervals
\(I\subset[P,2P]\),
\begin{equation}
 \sum_{p\in I}\chi(p)p^{it}
 =\delta_\chi\int_I\frac{u^{it}}{\log u}\,du
       +O((1+|t|)P^{\sigma_1}L^C).                  \label{eq:gr-primes}
\end{equation}
Here \(\delta_\chi=1\) for a principal character and zero otherwise,
and the estimate holds for every real \(t\). The exponent \(C\)
and implied constant may depend on \(a,B,B_1,\sigma_1\) and the
fixed bank modulus \(L_0\).
For progression banks the main vanishes unless \(\chi\) is induced
from a character modulo \(L_0\).
Removing modulus divisors or prime divisors of a nonzero target
\(|k|\le x^B\) preserves the error bound.
\end{lemma}

\begin{proof}
A fixed distance to the right of the zero-free boundary, the local
logarithmic-derivative formula and functional equation give
\begin{equation}
 \left|\frac{L'}{L}(s,\chi)+\frac{\delta_\chi}{s-1}\right|
                 \ll \log(q(|\Im s|+3)).
 \label{eq:gr-logderivative}
\end{equation}
Imprimitive Euler factors contribute \(O(\log q)\), since the real
part is fixed and positive.
Apply truncated Perron to
\(-L'/L(s,\chi)\), at
\(c=1+1/\log Z\), height \(Z^2\), and a half-integer endpoint
\(Z\in[P/2,4P]\).
The truncation error is \(O(\log^2 Z)\): near \(Z\), sum the bounds
\(O(Z\log Z/(Z^2|n-Z|))\); away from \(Z\), use absolute convergence
at \(c\).  Shift to \(\sigma_1\).
The principal pole gives \(\delta_\chi Z\), the new vertical integral
is \(O(Z^{\sigma_1}\log^2(qZ))\), and horizontal edges cost
\(O(Z^c\log(qZ)/(Z^2\log Z))\).
Prime powers cost \(O(Z^{1/2}\log^2 Z)\).
Replacing a real endpoint by the adjacent half-integer changes its
prime-weighted count by \(O(\log Z)\).  Therefore
\[
       \sum_{p\le Z}\chi(p)\log p
          =\delta_\chi Z+O(Z^{\sigma_1}L^C).
\]
Here \(Z\asymp P\le x^{B_1}\) and \(q\le x^B\) give
\(\log(qZ)\ll_{B,B_1}L\), justifying the displayed logarithmic cost.
Partial summation against \(u^{it}/\log u\) proves
\eqref{eq:gr-primes} with its linear \(1+|t|\) growth.

Expand the progression condition into the finitely many characters
modulo \(L_0\).  A constituent is principal exactly when the original
character is induced from \(L_0\).
There are at most \(O(\log(q(2+|k|)))=O(L)\) deleted prime divisors;
their absolute unweighted contribution is absorbed by the error.
\end{proof}

Every prime scale used below lies in one of the two bank ranges,
so \(P\le\exp(L^{2/5})\le x\) for large \(x\). Thus \(B_1=1\)
suffices for both comparison geometries; in the reflected application
the conductors are at most \(xL^{O(1)}\), and in the fixed-shift
application they are polynomial logarithmic.

In particular, for a nonexceptional character the prime factor saves
\(P^{\sigma_1-1}\) against its trivial length.  The all-height form
above is convenient: smooth product-cutoff Mellin transforms have
integrable \((1+|t|)\)-weighted absolute value bounded by a fixed
power of \(L\), so there is no boundary-smoothing loss from a sharp bank.


\subsection{Smooth prime sums at polynomial height}


The high-frequency argument uses a smooth version of
Lemma~\ref{lem:gr-primes}. Its height dependence is essential:
the linear \(1+|t|\) loss in the sharp-interval estimate is not
adequate when \(|t|\) can be of order \(x\).
Fix \(\sigma_1\in(11/12,1)\). For \(w\) smooth and supported on
\([1/2,2]\), with each required derivative bounded by a fixed
power of \(L\), the same zero-free input gives
\begin{equation}\label{eq:fixed-smooth-primes}
 \sum_p\chi(p)w(p/P)p^{it}
  =\delta_\chi\int_0^\infty\frac{w(u/P)u^{it}}{\log u}\,du
                       +O(P^{\sigma_1}L^{C_w})
\end{equation}
uniformly for \(q\le L^B\), characters \(\chi\bmod q\),
\(|t|\le x^C\), and \(P\ge\exp(cL^{1/8})\).
Here \(C_w\) depends on the fixed derivative bounds and on \(B,C\).
This is the smooth prime estimate used in
\cite[Lemma 5.1 and Section 6.4]{walsh}.

To obtain it directly from the present input, insert the von
Mangoldt weight and the smooth factor \(w(u/P)/\log u\), apply
Mellin inversion, and move the line from \(\Re s>1\) to
\(\Re s=\sigma_1\). Formula \eqref{eq:gr-logderivative} bounds
the regularized logarithmic derivative on this line by
\(O(\log(q(|\Im s|+3)))\). The Mellin transform is translated
in height by \(t\); its decay depends on the derivatives of
\(w\), rather than on \(|t|\). Its weighted integral is consequently
\(O(L^{C_w})\) even for \(|t|\le x^C\). The principal pole
gives the integral in \eqref{eq:fixed-smooth-primes}; prime powers
cost \(O(P^{1/2}L^{C_w})\), absorbed by the displayed error.
Truncation followed by dominated convergence justifies the contour
move. This uses the same fixed zero-free half-plane, without a
second prime-distribution theorem.

Impose simultaneously the bank class \(c\pmod{L_0}\) and a
prime class modulo \(q\). Incompatible classes contribute zero;
otherwise character orthogonality at \(\lcm(q,L_0)\) applies.
For large \(x\), every prime in the small bank exceeds
\(qL_0|h|\), so nonreduced prime classes and primes dividing
\(h\) do not occur. Integration by parts in the principal
integral, for every fixed \(j\), shows that the normalized
prime sum \(P^{-1}\sum_p\cdots\) is bounded by
\begin{equation}\label{eq:fixed-high-prime}
        L^{C_{w,j}}(1+|t|)^{-j}
                       +L^{C_w}\exp(-cL^{1/8}).
\end{equation}
All character and congruence costs here are polynomial logarithmic.

